\documentclass[10pt,a4paper]{amsart}
\usepackage[margin=19mm]{geometry}
\usepackage{amsmath,amssymb,mathtools}
\usepackage[scr=rsfs,cal=euler]{mathalfa}
\usepackage{xfrac}
\usepackage[unicode,hidelinks,bookmarks=false]{hyperref}
\newcommand{\ie}{i.e.\ }
\newcommand{\cf}{cf.\ }
\newcommand{\resp}{respectively\ }
\newcommand{\Subsection}{\subsection}
\providecommand{\nobreakdash}{\nobreak\hbox{-}}
\setlength{\emergencystretch}{2em}
\multlinegap0pt
\usepackage{bm}
\usepackage{bbm}
\usepackage{dsfont}
\usepackage{xspace}
\usepackage{cancel}
\usepackage{mathtools}
\usepackage{amsthm}
\makeatletter
\@ifundefined{theorem}{\newtheorem{theorem}{Theorem}}{}
\@ifundefined{lemma}{\newtheorem{lemma}[theorem]{Lemma}}{}
\makeatother
\newcommand{\Z}{\mathbb{Z}}
\title[On the distribution of shapes of sextic pure number fields]{On the distribution of shapes of sextic pure number fields}
\author{Anuj Jakhar, Ravi Kalwaniya, Anwesh Ray, Bidisha Roy}
\date{Source: arXiv:2601.09411v1 (CC BY 4.0)}
\begin{document}
\maketitle

\noindent\emph{Source and licence.} On the distribution of shapes of sextic pure number fields. Version arXiv:2601.09411v1 (CC BY 4.0), distributed under the Creative Commons Attribution 4.0 International licence, as recorded in the arXiv licence metadata for this exact version and shown on the abstract page https://arxiv.org/abs/2601.09411v1. Full rights notice: licenses/arxiv-2601.09411-CC-BY-4.0.txt.\par\smallskip

\noindent\textit{Editorial scope.} Counted unit: the Lemma whose source label is \texttt{lem:Rn-asymp} in the article (source0.tex lines 2168--2182, source label \texttt{lem:Rn-asymp}), delivered together with the article's own complete original proof of it (source0.tex lines 2185--2308, one \texttt{proof} environment of 124 lines). No mathematics is written, repaired or re-derived: statement and proof are the article's own text line for line. Required context retained uncounted: none. Every definition, display and label of those lines is retained unchanged. Omitted from this package and not redistributed: 1--2167, 2183--2184, 2309--2795. Nothing inside a retained excerpt is omitted or altered: every retained body is delivered exactly as the article's lines are written, so no omission receipt of any kind accompanies this package. Citations: the retained excerpts carry no citation command at all, so no bibliography entry of the article is transcribed and no reference is redistributed. Reference adaptation: none was needed and none was made; every reference inside the delivered unit resolves inside it, so no unresolved reference or citation is produced. Printed slips kept literal: no printed word, symbol, hypothesis or number of the counted statement or of its proof was altered. Layout and class: the article's own class is \texttt{amsart} with packages including amsmath, geometry, amscd, amsxtra, amsthm, amssymb, stmaryrd, xr, mathrsfs, mathtools, enumerate, commath, comment, enumitem; the wrapper is the lane's pinned \texttt{amsart} editorial wrapper at 10pt on A4, which restates the theorem-like environments the retained text needs and adds the article's own macro definitions that the retained text needs (\texttt{\textbackslash Z}), with extra wrapper packages (bm, bbm, dsfont, xspace, cancel, mathtools). The delivered numbering is the unit's own. Assets: no retained span contains a figure environment, a graphics-inclusion command, a PDF-page inclusion, a table or a TikZ picture; no image file is copied into this package, the package declares no asset dependency and no image file is redistributed with it. Licence: version arXiv:2601.09411v1 of this article is distributed under the Creative Commons Attribution 4.0 International licence, as recorded in the arXiv licence metadata for this exact version and shown on the abstract page https://arxiv.org/abs/2601.09411v1. A full rights notice is delivered at licenses/arxiv-2601.09411-CC-BY-4.0.txt. Characters: every retained excerpt is pure ASCII, so the delivered engine, XeLaTeX with its default Latin Modern text font, reports no missing character.
\begin{lemma}\label{lem:Rn-asymp}
Let $n$ be a squarefree positive integer divisible by $6$ and all primes $\ell\leq R_3$. Then one has
\[
\begin{split}\#\mathcal{C}^n(N,\mathscr{R})
=&
\frac{25 N^{1/5}}{124416}\prod_{\substack{\ell|n\\\ell\neq 2,3}}\left(1-\frac{3}{\ell^2}+\frac{2}{\ell^3}\right)
\bigl(R_1^{1/15}-R_1'^{1/15}\bigr)
\bigl(R_2'^{-2/15}-R_2^{-2/15}\bigr) \\
\times &\sum_{\substack{a_2a_4\in [R_3', R_3]\\
    a_2a_4\text{ is squarefree}}}\frac{n_{i,j}(a_2,a_4)\prod_{\substack{\ell|a_2a_4\\ \ell\neq 2,3}}\left(\frac{\ell-1}{\ell+2}\right)}{a_2^{3/5} a_4}+
O\!\left(N^{1/10}\right),
\end{split}\]
where the implied constant depends on $\varepsilon$, $n$, and the fixed
parameters $R_i,R_i'$.
\end{lemma}

\begin{proof}
Let $\Lambda_n=(N(n)\mathbb{Z})^5\subset\mathbb{Z}^5$. Since $n$ is squarefree,
the condition that a tuple $a=(a_1,\dots,a_5)$ be $\ell$-carefree for every
prime $\ell\mid n$ depends only on the reduction of $a$ modulo $\ell^2$, and
hence only on its reduction modulo $N(n)$. Consequently, the set
$\Omega_n\subset(\mathbb{Z}/N(n)\mathbb{Z})^5$ parametrizes exactly the residue
classes modulo $\Lambda_n$ that are allowed by the local carefree conditions such that $m$ is of Type $(i,j)$.
Fix once and for all a complete set of representatives
\[
\widetilde{\Omega}_n\subset ([0,N(n))\cap\mathbb{Z})^5
\]
for $\Omega_n$.

Every element of $\mathcal{C}^n(N,\mathscr{R})$ lies in a unique coset of $\Lambda_n$
corresponding to some $z\in\widetilde{\Omega}_n$, and therefore
\[
\mathcal{C}^n(N,\mathscr{R})
=
\bigsqcup_{z\in\widetilde{\Omega}_n}
\bigl((z+\Lambda_n)\cap\mathcal{C}(N,\mathscr{R})\bigr).
\]
Taking cardinalities and translating by $-z$, we obtain
\[
\#\mathcal{C}^n(N,\mathscr{R})
=
\sum_{z\in\widetilde{\Omega}_n}
\#\Bigl(\Lambda_n\cap(\mathcal{C}(N,\mathscr{R})-z)\Bigr)
=
\sum_{z\in\widetilde{\Omega}_n}
\#\Bigl(\mathbb{Z}^5\cap\bigl(\tfrac{1}{N(n)}\mathcal{C}(N,\mathscr{R})-\tfrac{z}{N(n)}\bigr)\Bigr).
\]

Fix $z=(z_1,\dots,z_5)\in\widetilde{\Omega}_n$. Writing
$a_i=z_i+N(n) b_i$ for $i=1,\dots,5$, the defining conditions of
$\mathcal{C}(N,\mathscr{R})$ translate into conditions on the integer vector
$b=(b_1,\dots,b_5)$ describing the shifted and rescaled region
\[
\tfrac{1}{N(n)}\mathcal{C}(N,\mathscr{R})-\tfrac{z}{N(n)}\subset\mathbb{R}^5.
\]
We now fix integers $b_2,b_4$ such that
\[
a_2=z_2+N(n) b_2>0,\qquad a_4=z_4+N(n) b_4>0,
\qquad\text{and}\qquad
a_2 a_4\in [R_3',R_3].
\]
For such a fixed pair $(a_2,a_4)$, we find that $(b_1,b_3,b_5)\in\mathbb{Z}^3$ must lie in the translated region
\[
\mathcal{M}\!\left(
\frac{N}{n^{26}a_2^4 a_4^4},
L_1',L_1,L_2',L_2
\right)
-
\frac{1}{N(n)}(z_1,z_3,z_5),
\]
where the shape parameters are given by
\[
L_1'=\sqrt{\frac{R_1' a_2}{a_4}},\qquad
L_1=\sqrt{\frac{R_1 a_2}{a_4}},\qquad
L_2'=n^6\frac{R_2' a_4}{a_2},\qquad
L_2=n^6\frac{R_2 a_4}{a_2}.
\]
We find that 
\[\begin{split}
&\#\Bigl(\mathbb{Z}^5\cap\bigl(\tfrac{1}{N(n)}\mathcal{C}(N,\mathscr{R})-\tfrac{z}{N(n)}\bigr)\Bigr)\\
=&\sum_{\substack{b_2,b_4\geq 0\\
(z_2+N(n)b_2)(z_4+N(n)b_4)\in [R_3', R_3]}}\# \left(\Z^3\cap\left(\mathcal{M}\!\left(
\frac{N}{n^{26}(z_2+N(n)b_2)^4 (z_4+N(n)b_4)^4},
L_1',L_1,L_2',L_2
\right)
-
\frac{1}{N(n)}(z_1,z_3,z_5)\right)\right),
\end{split}\]
where the values of $L_i$ and $L_i'$ are given above. We have that 
\[\begin{split}&\# \left(\Z^3\cap\left(\mathcal{M}\!\left(
\frac{N}{n^{26}(z_2+N(n)b_2)^4 (z_4+N(n)b_4)^4},
L_1',L_1,L_2',L_2
\right)
-
\frac{1}{N(n)}(z_1,z_3,z_5)\right)\right)\\
=& V\left(\mathcal{M}\!\left(
\frac{N}{n^{26}(z_2+N(n)b_2)^4 (z_4+N(n)b_4)^4},
L_1',L_1,L_2',L_2
\right)\right)+O(N^{\frac{1}{10}})\\
=& \frac{75}{8}\,\frac{N^{1/5}}{n^{\frac{26}{5}}(z_2+N(n)b_2)^{\frac{4}{5}} (z_4+N(n)b_4)^{\frac{4}{5}}}
\bigl(L_1^{2/15}-L_1'^{2/15}\bigr)
\bigl(L_2'^{-2/15}-L_2^{-2/15}\bigr)+O\left( N^{\frac{1}{10}}\right)\\
=&\frac{75}{8}\,\frac{N^{1/5}}{n^{6}(z_2+N(n)b_2)^{3/5} (z_4+N(n)b_4)}
\bigl(R_1^{1/15}-R_1'^{1/15}\bigr)
\bigl(R_2'^{-2/15}-R_2^{-2/15}\bigr)+O\left( N^{\frac{1}{10}}\right).
\end{split}\]
We deduce that 
\[\begin{split}&\#\mathcal{C}^n(N,\mathscr{R})\\
=&\frac{75}{8}\,\frac{N^{1/5}}{n^{6}}
\bigl(R_1^{1/15}-R_1'^{1/15}\bigr)
\bigl(R_2'^{-2/15}-R_2^{-2/15}\bigr) \sum_{z\in \widetilde{\Omega}_n}\sum_{\substack{b_2,b_4\geq 0\\
(z_2+N(n)b_2)(z_4+N(n)b_4)\in [R_3', R_3]}}\frac{1}{(z_2+N(n)b_2)^{3/5} (z_4+N(n)b_4)}\\
=&\frac{75}{8}\,\frac{N^{1/5}}{n^{6}}
\bigl(R_1^{1/15}-R_1'^{1/15}\bigr)
\bigl(R_2'^{-2/15}-R_2^{-2/15}\bigr) \sum_{a_2a_4\in [R_3', R_3]}\frac{\#\{z\in \widetilde{\Omega}_n\mid z_i\equiv a_i\pmod{N(n)}\text{ for }i=2,4\}}{a_2^{3/5} a_4}.
\end{split}\]
From the decomposition $\widetilde{\Omega}_n=\prod_{\ell|n} \widetilde{\Omega}_\ell$ one finds that 
\[\#\{z\in \widetilde{\Omega}_n\mid z_i\equiv a_i\pmod{N(n)}\text{ for }i=2,4\}=\prod_{\ell|n} \#\{z\in \widetilde{\Omega}_\ell\mid z_i\equiv a_i\pmod{\ell^2}\text{ for }i=2,4\}.\]
\par Let's count the number of $z\in \widetilde{\Omega}_\ell$ for which $z_2$ and $z_4$ are fixed. Assume that $\ell\neq 2,3$. If $z_2$ or $z_4$ is $0$, or if $\ell$ divides both $z_2$ and $z_4$ then there are no choices of $z$. Consider the case when $\ell\nmid z_2z_4$. In this case the number of choices for $(z_1, z_3, z_5)$ is $\ell^{6}\left(1-\frac{1}{\ell}\right)^2\left(1+\frac{2}{\ell}\right)$. Finally, consider the case when $\ell|z_2z_4$, but $z_2z_4\neq 0$. In this case none of the coordinates $(z_1, z_3, z_5)$ are divisible by $\ell$ and thus the number of choices in this case is $(\ell^2-\ell)^3=\ell^6(1-\frac{1}{\ell})^3$. Thus we have the following cases:
\[\#\{z\in \widetilde{\Omega}_\ell\mid z_i\equiv a_i\pmod{\ell^2}\text{ for }i=2,4\}=\begin{cases}
    0&\text{ if }\ell^2|a_2a_4,\\
    \ell^{6}\left(1-\frac{1}{\ell}\right)^2\left(1+\frac{2}{\ell}\right)&\text{ if }\ell\nmid a_2a_4,\\
    \ell^6(1-\frac{1}{\ell})^3&\text{ if }\ell|a_2a_4\text{ and }\ell^2\nmid a_2a_4.
\end{cases}\]
Note that by assumption on $n$, all primes that divide $a_2a_4$ must also divide $n$. Thus we deduce that 
\[\begin{split}
    & \sum_{a_2a_4\in [R_3', R_3]}\frac{\#\{z\in \widetilde{\Omega}_n\mid z_i\equiv a_i\pmod{N(n)}\text{ for }i=2,4\}}{a_2^{3/5} a_4}\\
    =& \sum_{\substack{a_2a_4\in [R_3', R_3]\\
    a_2a_4\text{ is squarefree}}}\frac{n_{i,j}(a_2,a_4)\prod_{\substack{\ell|n, \ell\nmid a_2a_4\\ \ell\neq 2,3}}\ell^{6}\left(1-\frac{1}{\ell}\right)^2\left(1+\frac{2}{\ell}\right)\times \prod_{\substack{ \ell\mid a_2a_4\\ \ell\neq 2,3}}\ell^6(1-\frac{1}{\ell})^3}{a_2^{3/5} a_4}\\
    =& (n/6)^6\prod_{\substack{\ell|n\\\ell\neq 2,3}}\left(1-\frac{3}{\ell^2}+\frac{2}{\ell^3}\right)\sum_{\substack{a_2a_4\in [R_3', R_3]\\
    a_2a_4\text{ is squarefree}}}\frac{n_{i,j} (a_2,a_4)\prod_{\substack{\ell|a_2a_4\\\ell\neq 2,3}}\left(\frac{\ell-1}{\ell+2}\right)}{a_2^{3/5} a_4}
\end{split}.\]
Thus we have shown that 
\[\begin{split}\#\mathcal{C}^n(N,\mathscr{R})=&\frac{25 N^{1/5}}{124416}\prod_{\substack{\ell|n\\\ell\neq 2,3}}\left(1-\frac{3}{\ell^2}+\frac{2}{\ell^3}\right)
\bigl(R_1^{1/15}-R_1'^{1/15}\bigr)
\bigl(R_2'^{-2/15}-R_2^{-2/15}\bigr) \\
\times& \sum_{\substack{a_2a_4\in [R_3', R_3]\\
    a_2a_4\text{ is squarefree}}}\frac{n_{i,j}(a_2,a_4)\prod_{\substack{\ell|a_2a_4\\ \ell\neq 2,3}}\left(\frac{\ell-1}{\ell+2}\right)}{a_2^{3/5} a_4}+
O\!\left(N^{1/10}\right).\end{split}\]
\end{proof}

\end{document}
